\scanpage{131}

% L1–8
\uline{Prop} (Cor du Thm) Soit $\varepsilon=\pm1$. La cat. des modules orth. lisses \inserted{$(E,q)$} sur $S$ de rg $n=2\nu+1$ est équivalente à la cat. pro\deleted{\unsure{d}}duit des couples $(E_0,\deleted{\unsure{L}}q_0;\omega_0,\Delta,\varphi)$, où $(E_0,q_0,\omega_0)$ est un module spécial quadratique pour $\varepsilon$, \deleted{\unsure{et}} $\Delta$ un mod. inv., $\varphi:L^{\otimes2}\overset{\sim}{\longleftarrow}\underline O_S$ (NB le couple $(L,\varphi)$ revient à la donnée de un torseur sous $\mu_2$). Le foncteur établissant l'équivalence est

% L9
\diagram{\begin{tikzcd}[ampersand replacement=\&,column sep=small]
 {(E_0,\ldots)}\arrow[r,mapsto] \& {\bigl(E_0\otimes\Delta,\ q_0\otimes L:}
 \& {E_0\otimes\Delta}\arrow[r]\arrow[rr,bend right=15] \& {\Delta^{\otimes2}\deleted{\unsure{\otimes}}}\arrow[r,"{\overset{\varphi^{-1}}{\sim}}"] \& {\underline O_S\bigr)}
\end{tikzcd}}

% L10
Le foncteur quasi-inverse est

% L11
\[
 (E,q)\longmapsto\bigl(\Delta=\det E,\ \varphi=\delta'(q)^{-1},\ E_0=E\otimes\Delta^{-1},
\]

% L12
\diagram{\begin{tikzcd}[ampersand replacement=\&,column sep=small]
 {q_0=\varepsilon q\otimes\Delta^{-1}:E_0}\arrow[r]\arrow[rr,bend right=15] \& \Delta^{-2}\arrow[r,"{{}^{t}\varphi}"] \& {\mathcal O_S,}
\end{tikzcd}}

% L13–14
\[
 \begin{aligned}
 &\omega_0\in\Gamma\det E_0\deleted{=\Gamma(\det E)\otimes L^{-n}\simeq\underline O_S}\\
 &\text{déduit de }\det E_0\simeq
 \underbrace{(\det \deleted{\unsure{E}}\inserted{E})}_{\deleted{\det E}\,\Delta}\otimes
 \underbrace{\Delta^{-n}}_{\Delta^{\unsure{4|-1|2\nu|-4}}}\simeq\mathcal O_S
 \end{aligned}
\]

% L15
\uline{NB} $\mathrm{Aut}\,V(q)\simeq O(q)/\mu_2\simeq SO(q)$, il est lisse conn\inserted{, réductif},

% L16
\deleted{\unsure{NB}}\quad\deleted{\unsure{et $\Longrightarrow$}}

% L17–19
\textcircled{3} Structure\unsure{s} de similitudes (forme quadratique \inserted{lisse} donnée loc. mod. scalaire inv, ou encore forme lisse à valeurs dans un $L$

% L20
\[
 q:E\longrightarrow L
\]

% L21–24
Soit $O'(q)$\inserted{$\subset\mathrm{Gl}(E)\times\mathrm{Gl}(L)=\mathrm{Gl}(E)\times\mathbb G_m$,} le groupe de ses autom, \deleted{il en fait}\deleted{\unsure{en fait}} formé des couples $(u,\lambda)$ tels que $q(ux)\overset{\mathrm{idt}}{=}\lambda q(x)$ \deleted{$SO(q)$} $O'(q)\unsure{\text{---}|\longrightarrow}\mathrm{Gl}(\unsure{n|u|E})$ est injectif, $O'(q)$ contient $\mathbb G_m$ comme centre (l'hom $O'(q)\mathbin{-}\mathbb G_m$ induit

% lower left
% archive sheet number: 130
